\section*{Informazioni preliminari}
Questo documento presenta il progetto \textit{PetMatch}, un'applicazione web sviluppata dal nostro gruppo per il corso di Tecnologie Web presso l'Università degli Studi di Padova.\\
\textbf{Il sito web è disponibile al link:}
\begin{center}
\url{https://tecweb.studenti.math.unipd.it/PetMatch/acanazza}
\end{center}

\subsubsection*{Credenziali di accesso}
Il sito web consente la visualizzazione come ospite, senza necessità di registrazione. Tuttavia sono presenti \textbf{due aree riservate agli utenti registrati}: l'area utente e l'area amministratore. Per accedere a queste aree, è possibile utilizzare le seguenti credenziali di test:
\begin{center}
\begin{tabular}{|l|c|c|}
\hline
\rowcolor{headerbrown}

\color{white}\textbf{Area} &
\color{white}\textbf{Username} &
\color{white}\textbf{Password} \\
\rowcolor{cellyellow} 
\hline
Area utente & \texttt{user} & \texttt{user} \\
\rowcolor{cellyellow} 
\hline
Area amministratore & \texttt{admin} & \texttt{admin} \\
\hline
\end{tabular}
\end{center}
In entrambi gli account stati predisposti dei dati di esempio per dimostrare le funzionalità del sito.\\\\
Si nota che queste credenziali sono valide esclusivamente su richiesta esplicita per la consegna del progetto e sono pertanto create appositamente da database; le credenziali create alla registrazione sul sito richiedono di rispettare specifici vincoli di sintassi.
